\chapter{การสังเคราะห์งานวิจัยเชิงวิพากษ์ระดับ Peer-Review}

การสังเคราะห์งานวิจัยตามมาตรฐานสากล ต้องยกระดับจากการบรรยายสรุป (Descriptive Summary) สู่การประเมินเชิงวิพากษ์ (Critical Appraisal) เพื่อจำแนกข้อเท็จจริงออกจากข้อจำกัดของการทดลอง

\section{โครงสร้างตารางวิเคราะห์ 7 คอลัมน์}

กรอบการวิเคราะห์ 7 คอลัมน์ (7-Column Research Matrix) ประกอบด้วยตัวบ่งชี้เชิงลึกที่ครอบคลุมมิติการวิจัย:
\begin{enumerate}
    \item \textbf{Claim (ข้ออ้างหลัก):} ข้อสรุปที่ผู้วิจัยพยายามพิสูจน์
    \item \textbf{Method (ระเบียบวิธี):} การออกแบบการทดลอง จำนวนซ้ำ และเครื่องมือวัด
    \item \textbf{Evidence (หลักฐานเชิงประจักษ์):} ข้อมูลตัวเลข ค่าเฉลี่ย ส่วนเบี่ยงเบนมาตรฐาน หรือระดับนัยสำคัญ $p$-value
    \item \textbf{Limitations (ข้อจำกัด):} ตัวแปรแทรกซ้อนที่ไม่ได้ควบคุม
    \item \textbf{Authority (ความน่าเชื่อถือ):} ความเชี่ยวชาญของคณะผู้วิจัยและสถาบันต้นสังกัด
    \item \textbf{Recency (ความสดใหม่):} ความก้าวหน้าเมื่อเทียบกับเทคโนโลยีปัจจุบัน
    \item \textbf{Network (เครือข่ายการอ้างอิง):} การเชื่อมโยงกับทฤษฎีรากฐานก่อนหน้า
\end{enumerate}

\section{การสังเคราะห์ข้อตกลงร่วมและข้อขัดแย้ง}

ในการเขียนสังเคราะห์ นักวิจัยต้องจัดหมวดหมู่ข้อค้นพบเป็น 2 ส่วนชัดเจน:
\begin{itemize}
    \item \textbf{Agreements (ข้อตกลงร่วม):} ประเด็นที่การทดลองจากหลายสถาบันให้ผลลัพธ์สอดคล้องกันอย่างมีนัยสำคัญ
    \item \textbf{Disagreements (ข้อขัดแย้ง):} จุดที่ผลการทดลองขัดแย้งกันเนื่องจากความแตกต่างของเงื่อนไขการทดลองหรืออุปกรณ์ที่ใช้
\end{itemize}

การแจกแจงเช่นนี้จะนำไปสู่การอภิปรายผลในบทที่ 4 และ 5 ของรายงานวิจัยที่มีน้ำหนักทางวิชาการสูง
